\section{Game Logic and Server Authority}
\label{sec:authority}

% ── Mark scheme: 15 marks ─────────────────────────────────────────────────────

The central tenet of the system architecture is that the STM32 microcontroller is the sole referee and single source of truth.
The PC client is treated as an untrusted, potentially hostile terminal.
This section details the move-validation pipeline, provides a step-by-step structural proof that rejected requests cannot corrupt the game state, and explains the win/draw detection algorithms.

\subsection{The Move Validation Pipeline Evaluates Legality in Fixed Order}

Every move request submitted to the server passes through the five-stage validation pipeline inside \texttt{game\_process\_move()} (\texttt{game.c}).
The order of checks is deterministic and strictly enforced to ensure predictable error codes when an incoming command violates multiple rules simultaneously:

\begin{enumerate}
    \item \textbf{Syntax and Parameter Parsing}: The server verifies that the payload contains exactly one colon separating a single player character from an index string (\texttt{<player>:<index>}). If the delimiter is missing or misplaced, or if the player token is not \texttt{'X'} or \texttt{'O'}, the server rejects the command with \texttt{NACK:BADCMD}.
    \item \textbf{Game State Check}: The server checks \texttt{g\_state}. If the state is \texttt{IDLE} (before \texttt{NEW}) or \texttt{GAME\_OVER} (match finished), moves are illegal, and the server returns \texttt{NACK:NOGAME}.
    \item \textbf{Index Range Verification}: The server ensures every character following the colon is an ASCII digit and parses the integer using \texttt{strtol()}. If non-digit characters are present or if the integer falls outside the closed interval $[0, 8]$, the server returns \texttt{NACK:RANGE}.
    \item \textbf{Turn Enforcement}: The player token is compared against \texttt{g\_current\_player}. If Player~O attempts a move when it is Player~X's turn (or vice versa), the server returns \texttt{NACK:NOTURN}.
    \item \textbf{Cell Occupancy Check}: The server indexes into the board array (\texttt{g\_board[idx]}). If the cell is not empty (i.e. contains \texttt{'X'} or \texttt{'O'} rather than \texttt{'.'}), the server returns \texttt{NACK:OCCUPIED}.
\end{enumerate}

Listing~\ref{lst:moveval} demonstrates the implementation of this strict evaluation cascade.

\begin{lstlisting}[language=C, caption={Move validation pipeline in \texttt{game\_process\_move()} (\texttt{game.c}).}, label=lst:moveval]
void game_process_move(const char *payload, char *resp, size_t max) {
    if (!payload) { snprintf(resp, max, "NACK:BADCMD\n"); return; }
    const char *colon = strchr(payload, ':');
    if (!colon || colon != payload + 1) {
        snprintf(resp, max, "NACK:BADCMD\n"); return;
    }
    char player = payload[0];
    if (player != 'X' && player != 'O') {
        snprintf(resp, max, "NACK:BADCMD\n"); return;
    }
    if (g_state != GAME_STATE_IN_PROGRESS) {
        snprintf(resp, max, "NACK:NOGAME\n"); return;
    }
    char *endptr;
    long idx = strtol(colon + 1, &endptr, 10);
    if (*endptr != '\0' || idx < 0 || idx > 8) {
        snprintf(resp, max, "NACK:RANGE\n"); return;
    }
    if (player != g_current_player) {
        snprintf(resp, max, "NACK:NOTURN\n"); return;
    }
    if (g_board[idx] != '.') {
        snprintf(resp, max, "NACK:OCCUPIED\n"); return;
    }
    /* Mutation only occurs here: all checks have passed */
    g_board[idx] = player;
    /* ... evaluate outcome and switch turn ... */
}
\end{lstlisting}

\subsection{Why Rejected Moves Do Not Change the Game State}

The project brief requires a clear explanation showing that a rejected request cannot alter server state.
This property is established through the structural encapsulation and control flow of the firmware:

\begin{enumerate}
    \item \textbf{Encapsulation of State Variables}: The server state is stored strictly in four file-scope static variables in \texttt{game.c}:
    \begin{itemize}
        \item \texttt{static char g\_board[10]} (the 9-cell board array),
        \item \texttt{static char g\_current\_player} (active player mark),
        \item \texttt{static GameState g\_state} (FSM state enum),
        \item \texttt{static char g\_winner} (outcome indicator).
    \end{itemize}
    No external module can modify these variables directly; all mutations occur exclusively inside \texttt{game.c}.
    \item \textbf{Gatekeeping Before Mutation}: In \texttt{game\_process\_move()}, every validation check (syntax, FSM state, index range, turn, and occupancy) executes an immediate return on failure.
    As evident in Listing~\ref{lst:moveval}, the first write to internal memory (\texttt{g\_board[idx] = player}) occurs strictly \emph{after} all five guard conditions have evaluated to true, following the defensive pre-condition pattern.
    \item \textbf{Framing and Unknown Commands}: If an incoming serial frame is malformed, exceeds 64 bytes, or specifies an unrecognised verb (e.g. \texttt{FOOBAR}), the rejection occurs inside \texttt{protocol.c}.
    The command dispatcher never invokes any function in \texttt{game.c}, ensuring zero side effects on game variables.
    \item \textbf{Test Suite Verification}: This invariant is verified empirically by automated unit tests in \texttt{test\_server.py}:
    \begin{itemize}
        \item \texttt{test\_out\_of\_turn\_move}: Asserts that sending an out-of-turn move returns \texttt{NACK:NOTURN} and confirms via \texttt{STATE} that \texttt{g\_board} is unmodified.
        \item \texttt{test\_occupied\_cell}: Verifies that playing on an occupied cell returns \texttt{NACK:OCCUPIED} without overwriting the existing mark.
        \item \texttt{test\_out\_of\_range\_move}: Verifies that indices outside 0 to 8 return \texttt{NACK:RANGE} with board state invariant.
        \item \texttt{test\_unknown\_command}: Verifies that garbage input returns \texttt{NACK:BADCMD} while the board remains unchanged.
    \end{itemize}
\end{enumerate}

\subsection{Sequential Turn Inversion Enforces Alternating Play Without Server-Side Inference}

The active turn is initialized to Player~X whenever \texttt{NEW} is accepted.
The server enforces alternation strictly: Player~X and Player~O must take turns sequentially.
Because our protocol explicitly includes the player token in the command (\texttt{MOVE:p:i}), the server checks both whether the move is legal and whether the claiming player matches \texttt{g\_current\_player}.

If a client attempts to bypass turn order (for example, by sending \texttt{MOVE:X:1} immediately after \texttt{MOVE:X:0}), the server detects the mismatch at check~4, issues \texttt{NACK:NOTURN}, and does not advance the turn.
Only after a valid move has been written to the board does the server toggle the active player:
\begin{equation}
\texttt{g\_current\_player} = (\texttt{player} == \text{'X'}) \mathbin{?} \text{'O'} : \text{'X'}
\end{equation}
Simultaneously, the physical LEDs update in hardware: the Red LED illuminates for X, and the Green LED illuminates for O.

\subsection{Deterministic Lookup Table Evaluates Win and Draw Conditions in Constant Time}

After a legal move is written to \texttt{g\_board}, the server evaluates whether the game has reached a terminal outcome before updating the turn indicator.

\subsubsection{Win Detection}
There are exactly eight winning 3-in-a-row lines in Tic-Tac-Toe: three horizontal rows, three vertical columns, and two diagonals.
These lines are encoded as an immutable lookup table:
\begin{align*}
\text{Rows:}\quad & \{0, 1, 2\},\ \{3, 4, 5\},\ \{6, 7, 8\} \\
\text{Columns:}\quad & \{0, 3, 6\},\ \{1, 4, 7\},\ \{2, 5, 8\} \\
\text{Diagonals:}\quad & \{0, 4, 8\},\ \{2, 4, 6\}
\end{align*}
The helper \texttt{check\_win(player)} iterates through these eight triads.
If all three cells match the player who just moved, the game terminates immediately.
The server transitions to \texttt{GAME\_OVER}, latches \texttt{g\_winner = player}, and transmits \texttt{WIN:<player>\textbackslash n}.
In hardware, the winner's status LED enters a 5\,Hz blinking cycle.

\subsubsection{Draw Detection}
Draw detection is evaluated \emph{after} win checking.
If no player has won, the server scans the board for empty cells (\texttt{'.'}).
If all nine cells are occupied and no winning line exists, the server transitions to \texttt{GAME\_OVER}, sets \texttt{g\_winner = 'D'}, and transmits \texttt{DRAW\textbackslash n}.
Both Red and Green LEDs blink simultaneously at 5\,Hz.

\subsubsection{Refusal of Subsequent Moves}
Once \texttt{g\_state} transitions to \texttt{GAME\_OVER}, check~2 in \texttt{game\_process\_move()} intercepts any subsequent \texttt{MOVE} or \texttt{RESIGN} commands and returns \texttt{NACK:NOGAME}.
The game cannot resume until an explicit \texttt{NEW} command or a hardware button press on \texttt{SW0} resets the board.

\subsection{Server-Side Validation of Client Requests}

In naive client-server games, the client computes game logic locally and sends updated state to the server.
In our system, the client is treated as an untrusted terminal:
\begin{itemize}
    \item \textbf{No Local Rules}: The Python GUI has zero knowledge of Tic-Tac-Toe win lines or valid cell occupancy. It merely forwards raw user clicks and renders whatever string the STM32 returns.
    \item \textbf{Resilience Against Modified Clients}: If a client script attempts to bypass button locking, send invalid coordinates (e.g. index 99), or transmit consecutive moves as Player~X, the STM32 referee rejects every invalid command at the serial communication boundary.
    \item \textbf{Auditable State}: The host can query \texttt{STATE} at any time to verify that the displayed grid matches the server's internal memory byte-for-byte.
\end{itemize}
This design ensures that physical rule enforcement remains authoritative regardless of client anomalies or link interference.
